\documentclass[10pt, a4paper]{article}
\usepackage[left=0.2in, right=0.2in, top=0.2in, bottom=0.2in]{geometry}
\usepackage{times}
\usepackage{enumitem}
\usepackage{hyperref}
\usepackage{xcolor}
\usepackage{tabularx}
\usepackage{array}

% Remove page numbering
\pagestyle{empty}

% Define colors for links
\hypersetup{
    colorlinks=true,
    linkcolor=blue!70!black,
    urlcolor=blue!70!black,
}

% Custom section formatting with consistent separators
\usepackage{titlesec}
\titleformat{\section}{\large\bfseries}{}{0em}{}[\vspace{2pt}\hrule height 0.5pt\vspace{4pt}]
\titlespacing{\section}{0pt}{0.5pt}{0pt}

% Remove default indentation
\setlength{\parindent}{0pt}

% Custom itemize formatting for clean bullet points
\setlist[itemize]{
    leftmargin=15pt,
    topsep=2pt,
    itemsep=1pt,
    parsep=1pt,
    label=\textbullet
}

% Define a command for subsection headers (bold with colon)
\newcommand{\skillitem}[1]{\textbf{#1} : }

\begin{document}

\begin{center}
    {\huge\textbf{YOUR FULL NAME}}
    
    \vspace{2pt}
    
    (+XX) XXXXXXXXXX \textbar{} City, Country
    
    \vspace{2pt}
    
    \href{mailto:your.email@example.com}{your.email@example.com} \textbar{} \href{https://www.linkedin.com/in/yourprofile/}{www.linkedin.com/in/yourprofile/} \textbar{} \href{https://www.github.com/yourusername}{www.github.com/yourusername} \textbar{} \href{https://yourportfolio.com/}{Portfolio}
\end{center}

\vspace{-12pt}

\section*{Education}

\textbf{University Name} \textbar{} City, State \hfill \textbf{Agg. CGPA: X.XX}\\
Degree Name \& Specialization \hfill Month YYYY - Month YYYY\\
\textit{One-line academic highlight or achievement (e.g., perfect score in final semester, dean's list, etc.)}

\vspace{5pt}

\section*{Technical Skills}

\skillitem{Languages}Language1 (proficiency), Language2 (proficiency), Language3 (proficiency)

\skillitem{Tools \& Technology}Tool1, Tool2, Tool3, Tool4, Tool5, Tool6, Tool7, Tool8, Tool9, Tool10

\skillitem{Frameworks}Framework1, Framework2, Framework3, Framework4, Framework5, Framework6, Framework7

\skillitem{Technical Domains}Domain1, Domain2, Domain3, Domain4, Domain5, Domain6, Domain7, Domain8, Domain9

\skillitem{Data Analysis \& Visualization}Skill1, Skill2, Skill3, Skill4, Skill5, Skill6, Skill7

\skillitem{Databases \& Cloud Platforms}Platform1, Platform2, Platform3

\vspace{5pt}

\section*{Projects}

\textbf{Project Title One} \textbar{} Tech Stack \& Key Technologies Used \hfill \href{https://github.com/yourusername/project-one}{\textcolor{blue!70!black}{Github Link}}
\begin{itemize}
    \item Brief description of what the project does and what problem it solves. Mention the core functionality and user-facing features.
    \item Key technical decision or implementation detail — security, architecture, or design choice that adds value.
    \item Quantifiable result or metric achieved (e.g., accuracy percentage, performance improvement, dataset size handled).
\end{itemize}

\vspace{3pt}

\textbf{Project Title Two} \textbar{} Tech Stack \& Key Technologies Used \hfill \href{https://github.com/yourusername/project-two}{\textcolor{blue!70!black}{Github Link}}
\textbar{} {\href{https://your-live-demo-link.com/}{\textcolor{blue!70!black}{Live Link}}}
\begin{itemize}
    \item What data or problem was analyzed/solved, and what was the overall scope (e.g., rows of data, time range, domain).
    \item Description of the data pipeline, feature engineering, or technical approach taken to tackle the problem.
    \item Business insight or measurable impact uncovered (e.g., percentage improvement, cost reduction, performance uplift).
\end{itemize}

\vspace{3pt}

\textbf{Project Title Three} \textbar{} Tech Stack \& Key Technologies Used \hfill \href{https://github.com/yourusername/project-three}{\textcolor{blue!70!black}{Github Link}}
\begin{itemize}
    \item End-to-end pipeline or system description — what data/input it processes and what output it produces.
    \item Technical methodology — tools, algorithms, or techniques used for data preparation and analysis.
    \item Key insight discovered (e.g., a metric, a trend, a business-relevant finding with concrete numbers).
\end{itemize}

\vspace{3pt}

\textbf{Project Title Four} \textbar{} Tech Stack \& Key Technologies Used \hfill \href{https://github.com/yourusername/project-four}{\textcolor{blue!70!black}{Github Link}} 
\textbar{} {\href{https://your-live-demo-link.com/}{\textcolor{blue!70!black}{Live Link}}}
\begin{itemize}
    \item What the system does at a high level — data sources, integrations, or APIs it connects to.
    \item Technical optimization made to improve system performance, accuracy, or scalability.
    \item Deployment approach and how it ensures reliability, groundedness, or real-time capability.
\end{itemize}

\vspace{3pt}

\textbf{Project Title Five} \textbar{} Tech Stack \& Key Technologies Used \hfill \href{https://github.com/yourusername/project-five}{\textcolor{blue!70!black}{Github Link}}
\begin{itemize}
    \item System architecture overview — what components were built and how they interact or compete.
    \item Key algorithmic or prompting strategy implemented to achieve the project's goal.
    \item Evaluation methodology or scoring logic designed, and what future use case it enables.
\end{itemize}

\vspace{3pt}

\textbf{Project Title Six} \textbar{} Tech Stack \& Key Technologies Used \hfill \href{https://github.com/yourusername/project-six}{\textcolor{blue!70!black}{Github Link}}
\begin{itemize}
    \item End-to-end pipeline description — model, tools, and overall workflow used to achieve the output.
    \item Model fine-tuning or optimization technique applied (e.g., quantization, LoRA, pruning) and its benefit.
    \item Additional module or feature integrated to extend functionality and improve user experience.
\end{itemize}

\vspace{3pt}

\textbf{Project Title Seven} \textbar{} Tech Stack \& Key Technologies Used \hfill \href{https://github.com/yourusername/project-seven}{\textcolor{blue!70!black}{Github Link}}
\begin{itemize}
    \item What problem the tool/extension solves, what it automates, and the workflow it improves for users.
    \item Core algorithm or detection logic implemented, along with the trigger mechanism and output format.
\end{itemize}

\vspace{5pt}

\section*{Experience}

\textbf{Job Title} \textbar{} Company Name, City, State \hfill Month YYYY - Present
\begin{itemize}
    \item Led [project or initiative name], performing [type of analysis or work] on [scope of data/system] to achieve [goal].
    \item Engineered [pipeline/system/tool] in [language/framework] to [what it does], creating [outcome or deliverable].
    \item Deployed [product/model/solution] achieving [accuracy/performance metric or business impact].
\end{itemize}

\vspace{3pt}

\textbf{Job Title} \textbar{} Company Name, City, State \hfill Month YYYY - Month YYYY\\
\textit{Project: Project Name | Tools/Technologies Used}
\begin{itemize}
    \item Engineered [type of pipeline or system] using [language/tool] to perform [ETL/analysis/processing] on [data volume].
    \item Conducted [type of analysis] to identify [key patterns, factors, or insights] related to [domain].
    \item Synthesized findings into [deliverable type] that [impact on team, product, or strategy].
\end{itemize}

\vspace{3pt}

\textbf{Job Title} \textbar{} Company Name, City, State \hfill Month YYYY - Month YYYY\\
\textit{Project: Project Name | Tools/Technologies Used}
\begin{itemize}
    \item Analyzed [problem domain] to design and build [product/system] in [tool/platform], improving [metric or experience].
    \item Investigated [N] [systems/applications/processes] to identify and resolve [N+] critical issues, demonstrating [skill].
\end{itemize}

\vspace{5pt}

\section*{Extracurricular Activities}

\textbf{Role/Position} \textbar{} Organization Name, Location \hfill Month YYYY - Month YYYY
\begin{itemize}
    \item Brief description of responsibilities, contributions, or what was built/organized during this activity.
\end{itemize}

\vspace{3pt}

\textbf{Role/Position} \textbar{} Organization Name, Location \hfill Month YYYY - Present
\begin{itemize}
    \item Brief description of testing, evaluation, or participation role and what products/programs were involved.
\end{itemize}

\vspace{3pt}

\textbf{Program/Competition Name} \textbar{} Institution Name, Mode \hfill Month YYYY - Month YYYY
\begin{itemize}
    \item Brief description of program outcome, ranking, or skill demonstrated.
\end{itemize}

\vspace{3pt}

\textbf{Program Name} \textbar{} Institution/Organization Name, Mode \hfill Month YYYY - Month YYYY
\begin{itemize}
    \item Brief description of modules completed or skills developed (e.g., leadership, communication, strategy).
\end{itemize}

\vspace{3pt}

\textbf{Role/Program Name} \textbar{} Organization Name, Mode \hfill Month YYYY - Month YYYY
\begin{itemize}
    \item Brief description of what you were selected for and what the program/fellowship involved.
\end{itemize}

\vspace{3pt}

\textbf{Role/Position} \textbar{} Organization Name, Mode \hfill Month YYYY - Month YYYY
\begin{itemize}
    \item Brief description of regional/community role, outreach efforts, or areas of contribution.
\end{itemize}

\vspace{3pt}

\textbf{Program Name} \textbar{} Organization Name, Mode \hfill Month YYYY - Month YYYY
\begin{itemize}
    \item Brief description of professional development completed (e.g., business analytics, strategic thinking).
\end{itemize}

\vspace{3pt}

\textbf{Role/Program Name} \textbar{} Organization Name, Mode \hfill Month YYYY - Month YYYY
\begin{itemize}
    \item Brief description of ambassadorship, representation, or campus/community role.
\end{itemize}

\vspace{5pt}

\section{Achievements}
\begin{itemize}
    \item Certified in [Topic/Skill] by [Institution], with demonstrated expertise in [specific area].
    \item Achieved [position/rank] in [competition name] at [event/fest], showcasing [skills demonstrated].
    \item Secured [position/rank] in [competition name] at [event/fest], demonstrating [skills demonstrated].
\end{itemize}

\vspace{5pt}

\section*{Publications}

\textbf{Full Title of Your Research Paper} \textbar{} Author One, Author Two
\begin{itemize}
    \item \textbf{Abstract:} One to two sentences summarizing the core finding, method, or contribution of the paper.
    \item \href{https://arxiv.org/abs/XXXXXXXXXX}{\textcolor{blue!70!black}{arXiv Pre-print, YYYY}} % or replace with journal/conference link
\end{itemize}

\vspace{5pt}
    
\section*{Certifications}
\begin{itemize}
    \item Passed [Certification Name] Certification. \href{https://credential-link.com/your-credential-id}{\textcolor{blue!70!black}{[View Credential]}}
\end{itemize}

\vspace{3pt}

\end{document}